\begin{problem}[第30届全国中学生物理竞赛复赛][变折射率劈尖干涉]{7}
一斜劈形透明介质劈尖，尖角为 $\theta$，高为 $h$。今以尖角顶点为坐标原点，建立坐标系如图(a)所示；劈尖斜面实际上是由一系列微小台阶组成的，在图(a)中看来，每一个小台阶的前侧面与 $xz$ 平面平行，上表面与 $yz$ 平面平行。劈尖介质的折射率 $n$ 随 $x$ 而变化，$n(x) = 1 + bx$，其中常数 $b > 0$。一束波长为 $\lambda$ 的单色平行光沿 $x$ 轴正方向照射劈尖；劈尖后放置一薄凸透镜，在劈尖与薄凸透镜之间放一档板，在档板上刻有一系列与 $z$ 方向平行、沿 $y$ 方向排列的透光狭缝，如图(b)所示。入射光的波面（即与平行入射光线垂直的平面）、劈尖底面、档板平面都与 $x$ 轴垂直，透镜主光轴为 $x$ 轴。要求通过各狭缝的透射光彼此在透镜焦点处得到加强而形成亮纹。已知第一条狭缝位于 $y = 0$ 处；物和像之间各光线的光程相等。
\begin{subproblem}
求其余各狭缝的 $y$ 坐标；
\end{subproblem}
\begin{subproblem}
试说明各狭缝彼此等距排列能否仍然满足上述要求。
\end{subproblem}
\begin{figure}[htbp]
    \centering
    \begin{subfigure}{0.45\textwidth}
        \centering
        \includegraphics[width=\linewidth]{image/1.png}
        \caption*{(a)}
    \end{subfigure}
    \hspace{1cm}
    \begin{subfigure}{0.3\textwidth}
        \centering
        \includegraphics[width=\linewidth]{image/2.png}
        \caption*{(b)}
    \end{subfigure}
\end{figure}
\end{problem}